\documentclass[8pt,landscape,letterpaper]{extarticle}
\usepackage[utf8]{inputenc}
\usepackage[margin=0.3in]{geometry}
\usepackage{amsmath,amssymb}
\usepackage{multicol}
\usepackage{xcolor}
\usepackage{tcolorbox}
\usepackage{enumitem}
\usepackage{titlesec}
\usepackage{booktabs}

% Colors
\definecolor{headercolor}{RGB}{25,25,112}
\definecolor{formulacolor}{RGB}{0,51,102}
\definecolor{criticalcolor}{RGB}{220,20,60}

% Compact spacing
\setlength{\parindent}{0pt}
\setlength{\parskip}{0pt}
\setlist{nosep, leftmargin=10pt}
\titlespacing*{\section}{0pt}{4pt}{2pt}
\titlespacing*{\subsection}{0pt}{3pt}{1pt}

% Section formatting
\titleformat{\section}{\large\bfseries\color{headercolor}}{\thesection}{0.5em}{}
\titleformat{\subsection}{\normalsize\bfseries\color{formulacolor}}{\thesubsection}{0.5em}{}

% Custom boxes
\newtcolorbox{keybox}[1][]{
  colback=yellow!10,
  colframe=orange!80,
  boxrule=0.5pt,
  arc=2mm,
  left=2pt,
  right=2pt,
  top=2pt,
  bottom=2pt,
  #1
}

\begin{document}

\begin{center}
{\LARGE\bfseries\color{criticalcolor} ECO310 FINAL EXAM CHEAT SHEET}\\
\vspace{2pt}
{\small Choice Under Uncertainty + Markets Theory}
\end{center}

\begin{multicols}{3}
\setlength{\columnseprule}{0.4pt}

\section*{UNCERTAINTY}

\subsection*{Expected Value \& Utility}
\begin{keybox}
\textbf{Expected Value:} $EV(L) = \sum_i p_i x_i$

\textbf{Expected Utility:} $EU(L) = \sum_i p_i u(x_i)$

\textbf{Certainty Equivalent:} $u(CE) = EU(L)$

\textbf{Risk Premium:} $RP = EV - CE$
\end{keybox}

\subsection*{Risk Aversion}
\begin{itemize}
\item \textbf{Risk Averse:} $u''(x) < 0$ (concave)
  \begin{itemize}
  \item $CE < EV$, $RP > 0$
  \item Prefers certainty over gamble
  \end{itemize}
\item \textbf{Risk Neutral:} $u''(x) = 0$ (linear)
  \begin{itemize}
  \item $CE = EV$, $RP = 0$
  \end{itemize}
\item \textbf{Risk Loving:} $u''(x) > 0$ (convex)
  \begin{itemize}
  \item $CE > EV$, $RP < 0$
  \end{itemize}
\end{itemize}

\textbf{Key Property:} Risk averse $\iff u(EV) > EU$

\subsection*{Common Utility Functions}
\begin{tabular}{ll}
\toprule
\textbf{Function} & \textbf{Derivatives} \\
\midrule
$u(x) = \sqrt{x}$ & $u' = \frac{1}{2\sqrt{x}}$, $u'' = -\frac{1}{4x^{3/2}}$ \\
$u(x) = \ln(x)$ & $u' = \frac{1}{x}$, $u'' = -\frac{1}{x^2}$ \\
$u(x) = x^\alpha$ & $u' = \alpha x^{\alpha-1}$, $u'' = \alpha(\alpha-1)x^{\alpha-2}$ \\
\bottomrule
\end{tabular}

\textbf{Note:} For $u(x) = \sqrt{x}$ or $\ln(x)$: Risk averse (concave)

\subsection*{Insurance}
\textbf{Setup:} Wealth $W$, loss $L$ w/ prob $p$

\textbf{Without Insurance:}
\[
EU = (1-p) u(W) + p u(W-L)
\]

\textbf{With Insurance:} Premium $\pi$, coverage $K$
\[
EU = (1-p) u(W-\pi) + p u(W-\pi-L+K)
\]

\textbf{Actuarially Fair:} $\pi = pK$

\textbf{Key Result:} Risk-averse person buys full insurance ($K=L$) if actuarially fair

\subsection*{Portfolio Choice}
\textbf{Setup:} Wealth $W$, fraction $\alpha$ in risky asset

\textbf{Safe return:} $r_f$ (certain)

\textbf{Risky return:} $r_h$ w/ prob $p$, $r_l$ w/ prob $1-p$

\textbf{Final Wealth:}
\begin{align*}
W_h &= W[(1-\alpha)(1+r_f) + \alpha(1+r_h)] \\
W_l &= W[(1-\alpha)(1+r_f) + \alpha(1+r_l)]
\end{align*}

\textbf{Expected Utility:}
\[
EU(\alpha) = p \cdot u(W_h) + (1-p) \cdot u(W_l)
\]

\textbf{FOC:}
\begin{keybox}
\[
p \cdot u'(W_h) \cdot W(r_h - r_f) + (1-p) \cdot u'(W_l) \cdot W(r_l - r_f) = 0
\]
\end{keybox}

\textbf{Check:} If $\alpha^* > 1$ or $< 0$, may be corner solution

\columnbreak

\section*{COST FUNCTIONS}

\subsection*{Deriving Costs from Production}
\textbf{Step 1:} Start with $q = f(L)$

\textbf{Step 2:} Invert: solve for $L(q)$

\textbf{Step 3:} Multiply by wage: $C(q) = w \cdot L(q)$

\textbf{Example:} $q = \sqrt{L}$, $w = 4$
\begin{align*}
&q = \sqrt{L} \implies L = q^2 \\
&C(q) = 4q^2
\end{align*}

\subsection*{Cost Concepts}
\begin{keybox}
\textbf{Total Cost:} $C(q) = FC + VC(q)$

\textbf{Marginal Cost:} $MC(q) = C'(q)$

\textbf{Average Cost:} $AC(q) = \frac{C(q)}{q}$

\textbf{AVC:} $AVC(q) = \frac{VC(q)}{q}$
\end{keybox}

\textbf{Key Facts:}
\begin{itemize}
\item MC intersects AC at $\min AC$
\item If $MC < AC$, then AC decreasing
\item If $MC > AC$, then AC increasing
\end{itemize}

\section*{PERFECT COMPETITION}

\subsection*{Characteristics}
\begin{itemize}
\item Many firms, homogeneous product
\item Firms are \textbf{price takers}
\item Take price $P$ as given
\end{itemize}

\subsection*{Profit Maximization}
\begin{keybox}
\textbf{Profit:} $\Pi(q) = P \cdot q - C(q)$

\textbf{FOC:} $\boxed{P = MC(q)}$

Produce where price equals marginal cost!
\end{keybox}

\textbf{Shutdown Rule:} Produce if $P \geq \min AVC$

Otherwise $q = 0$ (shut down)

\subsection*{Long-Run Equilibrium}
Free entry/exit drives profit to zero:
\begin{keybox}
\[
\boxed{P = MC = \min AC}
\]
\end{keybox}

Firms produce at minimum efficient scale.

\textbf{Zero Economic Profit:} $\Pi = 0$ in long run

\section*{MONOPOLY}

\subsection*{Characteristics}
\begin{itemize}
\item Single firm (only seller)
\item Firm is \textbf{price setter}
\item Faces entire market demand
\end{itemize}

\subsection*{Demand \& Revenue}
\textbf{Demand:} $Q^D(P)$ or inverse $P^D(Q)$

\textbf{Revenue:} $R(Q) = P^D(Q) \cdot Q$

\textbf{Marginal Revenue:}
\begin{keybox}
\[
MR(Q) = P^D(Q) + Q \cdot \frac{dP^D}{dQ}
\]

Since demand slopes down: $\boxed{MR < P}$
\end{keybox}

\subsection*{MR for Linear Demand}
\begin{keybox}
If $P = a - bQ$, then $\boxed{MR = a - 2bQ}$

\textbf{Same intercept, TWICE the slope!}
\end{keybox}

\subsection*{Profit Maximization}
\begin{keybox}
\textbf{Profit:} $\Pi(Q) = P^D(Q) \cdot Q - C(Q)$

\textbf{FOC:} $\boxed{MR(Q) = MC(Q)}$
\end{keybox}

\textbf{Steps:}
\begin{enumerate}[leftmargin=10pt]
\item Convert to inverse demand: $P^D(Q)$
\item Find MR (for linear: double slope)
\item Find MC from $C(Q)$
\item Set $MR = MC$, solve for $Q^*$
\item Find price: $P^* = P^D(Q^*)$
\end{enumerate}

\columnbreak

\section*{COMPETITION vs MONOPOLY}

\begin{tabular}{lcc}
\toprule
\textbf{Feature} & \textbf{Competition} & \textbf{Monopoly} \\
\midrule
\# of firms & Many & One \\
Profit max & $P = MC$ & $MR = MC$ \\
Price vs MC & $P = MC$ & $P > MC$ \\
Output & Higher & Lower \\
Price & Lower & Higher \\
LR Profit & Zero & Positive \\
Efficiency & Yes & No (DWL) \\
\bottomrule
\end{tabular}

\section*{WELFARE}

\subsection*{Surplus Measures}
\textbf{Consumer Surplus (CS):}
\[
CS = \int_0^Q P^D(q) \, dq - P \cdot Q
\]
Area under demand, above price

\textbf{Producer Surplus (PS):}
\[
PS = P \cdot Q - \int_0^Q MC(q) \, dq
\]
Area above MC, below price

\textbf{Total Surplus:} $TS = CS + PS$

\textbf{Deadweight Loss:}
\[
DWL = TS_{\text{efficient}} - TS_{\text{actual}}
\]

\subsection*{Welfare Theorems}
\begin{keybox}
\textbf{First Welfare Theorem:}

Competitive equilibrium is Pareto efficient

\textbf{Second Welfare Theorem:}

Any Pareto efficient allocation achievable via redistribution + competition
\end{keybox}

\section*{GENERAL EQUILIBRIUM}

\subsection*{Pareto Efficiency}
\textbf{Definition:} Can't make someone better off without making someone else worse off

\textbf{Condition:}
\begin{keybox}
\[
\boxed{MRS_A = MRS_B}
\]
\end{keybox}

If $MRS_A \neq MRS_B$, gains from trade exist

\subsection*{Edgeworth Box}
\textbf{Setup:} 2 consumers, 2 goods

\textbf{Endowments:} $(\omega_A^X, \omega_A^Y)$, $(\omega_B^X, \omega_B^Y)$

\textbf{Contract Curve:} Set of Pareto efficient allocations

\textbf{Competitive Equilibrium:} Where budget lines tangent to indifference curves at same point

\section*{QUICK REFERENCE}

\subsection*{Key Derivatives}
\begin{align*}
\frac{d}{dx}(x^n) &= nx^{n-1} \\
\frac{d}{dx}(\ln x) &= \frac{1}{x} \\
\frac{d}{dx}(e^x) &= e^x \\
\frac{d}{dx}(\sqrt{x}) &= \frac{1}{2\sqrt{x}}
\end{align*}

\subsection*{Problem Type Recognition}
\textbf{``Expected value''} $\rightarrow$ $\sum p_i x_i$

\textbf{``Expected utility''} $\rightarrow$ $\sum p_i u(x_i)$

\textbf{``Certainty equivalent''} $\rightarrow$ Solve $u(CE) = EU$

\textbf{``Risk premium''} $\rightarrow$ $EV - CE$

\textbf{``Insurance''} $\rightarrow$ Compare $EU$ with/without

\textbf{``Portfolio''} $\rightarrow$ Maximize $EU(\alpha)$

\textbf{``Cost function''} $\rightarrow$ Invert production function

\textbf{``Competitive firm''} $\rightarrow$ $P = MC$

\textbf{``Monopolist''} $\rightarrow$ $MR = MC$

\textbf{``Long-run''} $\rightarrow$ $P = MC = \min AC$

\subsection*{Common Mistakes}
\begin{itemize}
\item Don't confuse EV and EU!
\item Must invert production function first
\item Monopoly uses $MR = MC$ not $P = MC$
\item For linear demand: MR has 2× slope
\item Convert demand to inverse for monopoly
\item Check corner solutions in portfolio
\end{itemize}

\end{multicols}

\begin{center}
\begin{tcolorbox}[colback=yellow!20,colframe=orange!80,width=0.95\textwidth,arc=3mm,boxrule=1pt]
\centering
\textbf{MASTER THESE 5 THINGS:} (1) Expected Utility vs EV, (2) Certainty Equivalent, (3) MR for Linear Demand, (4) $P=MC$ (competition) vs $MR=MC$ (monopoly), (5) Inverting Production Functions
\end{tcolorbox}
\end{center}

\end{document}
